\section{The incoming Pell growth arithmetic interface}\label{sec:incoming210}\label{g210:sec:interface}

This section retains the complete incoming Pell Growth 210 arithmetic
argument, rather than claiming it as a new result of this round~\cite{growth210}. Its source
and independent fresh replay are identified in Section~\ref{sec:comparison}.
The later Round23 affine horizon and coefficient changes compose with this
argument at explicitly checked rows.

We describe the semantic coordinates before the exact projections. They are
not all separate coordinates after projection; the incoming finite residual array is supplied in its source closure.

The positive horizon has the cofinal form
\begin{equation}\label{g210:eq:horizon}
t=16+4z^2.
\end{equation}
The coordinate $F$ provides formula weights $1+jF$. The quantity
$B=\binom{F+t}{t}$ will be recovered from a radix digit, and $N=B^x$ is the
generalized CRT modulus. The beta spacing is $y=F(2v+1)$ and the beta code
$x$ satisfies
\begin{equation}\label{g210:eq:betaend}
x=F+1+(y+1)r_0,
\qquad x=1+[1+(t+1)y]r_t.
\end{equation}
Eight independent CRT representatives encode the four local parameters, two
pending states, and two beta quotients. The ninth representative $C$ encodes
the row index. It is explicitly even and satisfies
\begin{equation}\label{g210:eq:selector}
B\eta=C+F+1.
\end{equation}

The shared Pell machinery has its own large center $A_*$ and a square frame
containing every actual requested base, exponent and output. Its inherited
signed index theorem first authenticates the principal Pell pairs at their
nominal, initially signed indices. Its distinguished anchor is the pair for
$C$, and its cap proves
\begin{equation}\label{g210:eq:cap}
C\ne0,\qquad |k|<|C|\quad\hbox{for every other retained nominal index }k.
\end{equation}
After the new growth pair is removed there are 13 retained principal pairs
and 15 exact-power requests. These include the requests at base $q+1$ with
indices $F+t$, all eight ordinary representatives, and $C$; the requests at
base $q$ with indices $t$, $y+1$, and $x$; and the two singleton requests
with exponent $x$ and bases $B$ and the radix base.

For any integer base $b$ with $|b|\ge2$, the inherited signed power theorem
forces a nonnegative requested exponent $k$ and the exact output $b^k$.
This theorem is used only after the main pairs have been oriented. The frame
and anchor theorem does not require the front coordinates to have recovered
their arithmetic signs. The new independent Pell witnesses will not be
members of this frame; their existing output coordinate $F$ remains a frame
atom.

The incoming source changes to this interface are checked by independently
reconstructing all 34 retained square-frame atoms, all 13 principal norms and
multiplication wires, the whole orientation sum and the anchor cap.

\subsection{A sharper bound for the same local syntax}

The source-bound constructor audit gives maximum depth 19 for a local formula
label. The nine axiom terms have depths
\[
4,3,3,3,3,5,4,3,19.
\]
Every leaf parameter is strictly below $t$. The following sharper majorant
uses the same expression trees; it does not change the calculus.

\begin{lemma}[A capacity bound for constructor trees]\label{g210:lem:capacity}
Put $d=2^{19}$ and $L=8^{d-1}t^d$. Every local formula label is strictly
less than $L$.
\end{lemma}
\begin{proof}
For $0\le u,v\le m$, the largest constructor is bounded by $E(m,m)$, and
\[
E(m,m)+1=8m^2+14m+8\le8(m+1)^2.
\]
A depth-zero leaf has value plus one at most $t$. The recurrence
$H_0=t$, $H_{r+1}=8H_r^2$ therefore bounds value plus one at depth $r$.
Its explicit solution is $H_r=8^{2^r-1}t^{2^r}$. At depth 19 it equals
$L$. Unequal child depths are handled by monotonicity. Thus a label has value
at most $L-1$, as asserted.
\end{proof}

\section{Growth from three quadratic equations}\label{g210:sec:growth}

The old exact request $F=t^{2^{20}t}$ is stronger than the proof extraction
needs. The required interface is a sufficiently large even nonzero $F$,
eventually positive, with a cofinal completeness construction. Introduce
three integer coordinates $X,Z,u$ and impose
\begin{equation}\label{g210:eq:growth}
\boxed{\quad Z=tF,\qquad X^2-16Z^2-8ZF=1,\qquad
F=2^{19}t(2u+1).\quad}
\end{equation}
These are quadratic equations. They replace five coordinates of the old
growth request: its principal Pell triple, its modulus and its quotient.
The obsolete growth index is removed from the frame. The old principal pair's
norm, multiplication wire, orientation term and cap term are removed exactly
as specified by the source certificate; $F$ remains a frame atom.

\subsection{Pell recurrence facts}

For $a\ge2$ define $X_n(a),Y_n(a)$ by
\begin{align*}
X_0&=1,&Y_0&=0,\\
X_{n+1}&=aX_n+(a^2-1)Y_n,&Y_{n+1}&=X_n+aY_n.
\end{align*}
The least positive Pell unit at discriminant $a^2-1$ is
$a+\sqrt{a^2-1}$: the ordinate of a nontrivial positive integral solution
is at least one, and $(a,1)$ is a solution. The standard Pell classification
therefore gives, for every integer solution with $Y\ne0$,
$|X|=X_n(a)$ and $|Y|=Y_n(a)$ at a unique $n\ge1$.

We need three elementary consequences:
\begin{align}
Y_n(a)&\equiv n\pmod{a-1},\label{g210:eq:indexcong}\\
\vp(Y_n(a))&=\vp(n) &&(a\hbox{ odd},\ n\ge1),\label{g210:eq:valuation}\\
Y_n(a)&\ge(2a-1)^{n-1} &&(n\ge1).\label{g210:eq:pellgrowth}
\end{align}
The first follows by specializing the integer polynomial recurrence at
$a=1$. For the second, all $X_n(a)$ are odd and $Y_n(a)$ has the parity of
$n$ when $a$ is odd; the duplication identity
$Y_{2m}=2X_mY_m$ adds exactly one factor of two at each doubling. For the
third, $Y_n$ is increasing and
$Y_{n+1}=2aY_n-Y_{n-1}\ge(2a-1)Y_n$.

\begin{theorem}[Sign-free growth]\label{g210:thm:growth}
Every integer solution of \eqref{g210:eq:horizon} and \eqref{g210:eq:growth} has even
nonzero $F$ and
\[
|F|\ge L^t,\qquad L=8^{2^{19}-1}t^{2^{19}}.
\]
In particular $|F|\ge t^t>t$.
\end{theorem}
\begin{proof}
The last equation makes $F$ nonzero. Eliminating $Z$ gives
\[
X^2-((4t+1)^2-1)F^2=1.
\]
Pell classification at the odd center $a=4t+1$ gives
$|F|=Y_n(a)$ for $n\ge1$. Put $d=2^{19}$. Since $4t\mid F$,
\eqref{g210:eq:indexcong} gives $4t\mid n$. Moreover,
\[
\vp(n)=\vp(F)=19+\vp(t).
\]
Consequently the positive integer $n/(4t)$ has exact valuation 17, and
$n\ge dt$. Formula \eqref{g210:eq:pellgrowth} gives
\[
|F|\ge(8t+1)^{n-1}\ge(8t)^{dt-1}\ge
\bigl(8^{d-1}t^d\bigr)^t.
\]
The last ratio is exactly $8^{t-1}/t\ge1$. This proves the result with no
assumption about the sign of $F$.
\end{proof}

The flexibility of this certificate is useful. It does not insist that $F$
be one specified power. Section~\ref{g210:sec:cofinal} will choose among its Pell
ordinates to meet an additional binary-residue condition. Without that
choice, merely proving a growth lower bound would not establish completeness.

\section{The dominant CRT index supplies the radix}\label{g210:sec:radix}

Replace the old base $B[2+4(C+\sum A_i)]$ by
\begin{equation}\label{g210:eq:newradix}
D=2(C+F+1),\qquad q=D^x.
\end{equation}
The power in this display is still represented by the retained polynomial
certificate. It is not an operation in the final polynomial. The definition
of $D$ is affine, so its old multiplication coordinate and defining residual
can be eliminated without raising degree. Crucially, the actual affine
expression is substituted into the main square frame as well.

Equation \eqref{g210:eq:selector} gives $B\mid D$ directly. Once the signs have
been established, the existing anchor supplies $D>2C$ and bounds every
needed binomial upper argument by $C$. These two roles replace the old
explicit factor $B$ and the old sum-majorant.

\subsection{Recovering the signs in the correct order}\label{g210:sec:signs}

The main shared-Pell frame and $C$-anchor theorem apply first, while the front
coordinates are signed. Theorem~\ref{g210:thm:growth} independently gives even
nonzero $F$ with $|F|>t$. Because $C$ is even, $C+F+1$ is odd and nonzero;
thus $|D|\ge2$ and the request $q=D^x$ forces $x\ge0$.

Since $y=F(2v+1)$, $|y|\ge|F|$. If $x=0$, the terminal endpoint in
\eqref{g210:eq:betaend} would make $1+(t+1)y$ divide $-1$. This is impossible:
\[
|1+(t+1)y|\ge(t+1)|y|-1>1.
\]
Hence $x\ge1$ and $|q|\ge2$. The retained power at base $q$ and exponent
$y+1$ now forces $y+1\ge0$. Negative $y$ would be at most $-|F|<-1$.
Therefore $y\ge|F|>0$.

The value $x=1$ needs a separate exclusion. The first endpoint would give
$-F=(y+1)r_0$, whereas $0<|F|\le y<y+1$. This is impossible. Thus
$x\ge2$, and the terminal endpoint gives
\begin{equation}\label{g210:eq:xbound}
x\ge2+(t+1)y.
\end{equation}
It follows that $|q|\ge2^x\ge4$ and $|q+1|\ge3$. The base-$(q+1)$
requests now force all actual CRT representatives and $F+t$ to be
nonnegative. If $F$ were negative, then $F+t\le t-|F|<0$; hence $F>0$.
The already nonzero $C$ is positive as well. Thus $D>0$ and $q>0$.

Finally, if $r_0\le-1$, the first endpoint gives $x\le F-y\le0$.
Therefore $r_0\ge0$ and $x\ge F+1$. We have recovered every required sign
without adding a sign witness. In particular, the old assumption of a positive
growth output has been replaced by a proved consequence of retained equations.

\subsection{One binomial bound for the whole front}

By \eqref{g210:eq:cap}, the initial upper argument $F+t$ and all ordinary
representatives are below $C$; the anchor argument itself is $C$.
The lower cutoffs $t,y+1,x$ are at most $x$. Also
\[
D=2(C+F+1)>2C,\qquad q=D^x>(2C)^x.
\]
For every binomial coefficient in any required low digit, including $B$,
\begin{equation}\label{g210:eq:uniformbin}
0\le\binom{a}{j}\le C^x<q/2^x
\quad(0\le a\le C,\ 0\le j\le k\le x).
\end{equation}
Coefficients with $j>a$ are zero; $j=0$ is included since $C>t\ge16$.
As $x\ge F+1$, this implies $C^x/(q-1)<1/32$.

For clarity, the retained remainder certificate uses $M=q^x$, denominators
$q^k$, integer ratios $a_k=M/q^k$, and normalized remainders $\rho_k$.
Its exact division equations imply $a_k\mid\rho_k$. Its ternary norm budget
\[
M^2=1+4\sum_k\rho_k^2+r^2+s^2+w(w+1)
\]
gives $|\rho_k|<M/2$, hence the actual integer remainder
$r_k=\rho_k/a_k$ satisfies $|r_k|<q^k/2$. The genuine binomial low tail is
less than $q^k/32$ by \eqref{g210:eq:uniformbin}. Congruence modulo $q^k$
therefore forces the two remainders to agree exactly. This is an equality
argument, not an approximation.

The first digit consequently gives $B\equiv\binom{F+t}{t}\pmod q$.
Equation \eqref{g210:eq:selector}, with nonzero integer $\eta$, also gives
\[
|B|\le C+F+1=D/2<q/2.
\]
The true binomial coefficient is positive and below $q/32$, so
\begin{equation}\label{g210:eq:Bexact}
B=\binom{F+t}{t}>0.
\end{equation}
It is at least two. The retained request authenticates $N=B^x$, and
\[
q=D^x=(2B\eta)^x=N(2\eta)^x.
\]
Thus $N\mid q$, the exact modulus relation required by the generalized
range compiler. The same tail argument authenticates all remaining range
digits. The inherited local arithmetic bound below $2^x$ continues to hold:
it needs $F\ge t^t$, $t\ge16$ and \eqref{g210:eq:xbound}, all now proved.

\section{A cofinal choice with an odd binomial selector}\label{g210:sec:cofinal}

Soundness alone is not enough. The old exact power supplied high powers of
two automatically. The new ordinate has a prescribed finite valuation in our
completeness construction, so we must choose its other binary digits.

\begin{lemma}[Pell permutation modulo powers of two]\label{g210:lem:permutation}
If $a\equiv1\pmod4$, the map $s\mapsto Y_s(a)$ induces a permutation of
$\Z/2^r\Z$ for every $r\ge1$. Each odd residue has arbitrarily large
positive odd indices realizing it.
\end{lemma}
\begin{proof}
All $X_s(a)\equiv1\pmod4$ and $Y_s(a)\equiv s\pmod2$. Put $h=2^r$.
Equation \eqref{g210:eq:valuation} gives $\vp(Y_h)=r$. The norm and
$8\mid a^2-1$ imply $\vp(X_h^2-1)\ge2r+3$. Since $X_h+1$ has valuation
one, $X_h\equiv1\pmod{2^{2r+2}}$. The addition identity then gives
\[
Y_{s+h}-Y_s=(X_h-1)Y_s+X_sY_h\equiv2^r\pmod{2^{r+1}}.
\]
The two index lifts $s$ and $s+2^r$ produce the two distinct output lifts.
Induction from modulo two proves bijectivity and periodicity. Positive lifts
in the resulting residue class are cofinal and preserve odd parity.
\end{proof}

\begin{theorem}[Compatible cofinal growth witnesses]\label{g210:thm:cofinal}
There are cofinally many horizons of the form \eqref{g210:eq:horizon} for which
\eqref{g210:eq:growth} has arbitrarily large positive solutions $F$ such that
$\binom{F+t}{t}$ is odd.
\end{theorem}
\begin{proof}
Choose the free horizon parameter $z$ divisible by $2^{12}$. These horizons
remain cofinal and satisfy $t\equiv16\pmod{2^{26}}$. Thus $\vp(t)=4$ and
the binary digit of $t$ at position 23 is zero.

Put $d=2^{19}$, $n_0=dt$, $y_0=Y_{n_0}(4t+1)$ and
$a_0=X_{n_0}(4t+1)$. The congruence modulo $4t$ and
\eqref{g210:eq:valuation} show that $b_0=y_0/(dt)$ is an odd integer.
Moreover $a_0\equiv1\pmod4$, since $n_0$ is even and
$X_{2m}=2X_m^2-1$ with $X_m$ odd.

Choose $\ell>26$ with $2^\ell>t$, and set $r=\ell-23$. The odd factors
$t/16$ and $b_0$ are units modulo $2^r$. By Lemma~\ref{g210:lem:permutation},
choose a positive odd $s$, as large as desired, such that
\[
(t/16)b_0Y_s(a_0)\equiv1\pmod{2^r}.
\]
Set $n=n_0s$ and
\[
F=Y_n(4t+1)=y_0Y_s(a_0),\quad X=X_n(4t+1),\quad Z=tF,
\quad u=\frac{F/(dt)-1}{2}.
\]
Composition of Pell units proves the displayed identity for $F$, and the odd
quotient makes $u$ integral. These values satisfy all three equations and
\begin{equation}\label{g210:eq:binarytarget}
F\equiv2^{23}\pmod{2^\ell}.
\end{equation}
Since $t<2^\ell$ and its bit 23 is zero, the binary supports of $F$ and $t$
are disjoint. There are no carries in their addition. Writing $s_2(m)$ for
the binary digit sum, the elementary factorial identity gives
\[
\vp\binom{F+t}{t}=s_2(F)+s_2(t)-s_2(F+t)=0.
\]
Thus the binomial is odd. Cofinal positive lifts of $s$ give arbitrarily
large $F$ at each chosen horizon.
\end{proof}

The congruence proof is finite arithmetic at an arbitrarily chosen modulus;
no analytic inverse theorem or computation of an infinite 2-adic object is
needed. The retained executable modular examples check this construction
at the actual $d=2^{19}$, but the theorem follows from the recurrence proof
for every modulus.

\subsection{The beta moduli need only prime support}

The new construction has $\vp(F)=23$ at the selected horizons. We must not
reuse the old sufficient condition $t!\mid F(2v+1)$ for arbitrary large $t$.
Fortunately the beta CRT needs only every prime at most $t$ to divide $y$.

Choose a sufficiently large positive odd $w$ divisible by all odd primes at
most $t$, and put $y=Fw$. The integer $y$ is divisible by every prime at
most $t$. If a prime $p$ divided both $1+iy$ and $1+jy$ for
$1\le i<j\le t+1$, it would be coprime to $y$ and divide $j-i$.
Then $p\le t$, contradicting $p\mid y$. Thus the beta moduli are pairwise
coprime. The odd multiplier can also make $y$ exceed every finite pending
state and $F$, and it has exactly the required form $2v+1$.

This replacement of an unnecessarily strong factorial divisibility is part
of the completeness proof. Keeping the old assertion would leave a real gap.

\section{From a zero to a refutation, and back}\label{g210:sec:calibration}

\subsection{Sharp separation of the weighted labels}

Apply Lemma~\ref{lem:separate} to the sharper capacity bound.
The equality case $F=L^t$ is permitted by its strict coefficient-difference
argument.

The decoded local Horn rows have weighted updates
\[
S_{i+1}(1+Fv_i)=S_i(1+Fp_i)(1+Fq_i),
\]
with missing premises filled by the truth label zero. The endpoint $y=F$
would give an initial beta residue zero and terminal residue one. Positive
weights propagate a zero initial state to zero at every row, so this case
is impossible. Hence $y>F$, the actual initial state is $F+1$ and the final
state is one. All states are positive. Multiplying the updates and cancelling
states gives the equality between the conclusion multiset and the premise
multiset plus the target label one. Theorem~\ref{g210:thm:growth} and
Lemma~\ref{g210:lem:capacity} give precisely the hypotheses of
Lemma~\ref{lem:separate}.

Cycles in a selected multiset of Horn rules do not cause a gap. Let $U$ be
the labels not derivable from the available axioms using those rules. Every
selected rule whose conclusion belongs to $U$ has at least one premise in
$U$. If the target one belonged to $U$, multiplicity balance would say
that the number of conclusions in $U$ is one more than the number of
premises in $U$, contradicting the preceding inequality. Thus there is a
finite $\Ccal$ refutation.

\subsection{Soundness of the full arithmetic source}

From an arbitrary integer zero, first classify the standalone Pell pair to
obtain even, nonzero large $F$. Independently use the retained main frame,
anchor and orientation theorem to recover the canonical signed indices.
Section~\ref{g210:sec:signs} then proves the front signs. The uniform digit
argument recovers $B$, $N$, and all range binomials. The inherited generalized
binomial-CRT theorem applies because $F\ge L^t\ge t^t$ and hence
$F\ge\operatorname{lcm}(1,\ldots,t)^2$. Its two required conclusions are
retained: each row has a full prime-power selector exceeding $t$, and every
full prime power of $B$ is assigned to some row. The local arithmetic bound
below $2^x$ converts those congruences into actual local equations.

This decodes the beta states, valid local rules, and their strict parameter
ranges. The preceding product and Horn-flow argument yields a finite
$\Ccal$ refutation. All semantic source changes have now been accounted for;
the rest of the source decoding is the identified inherited theorem.

\subsection{Constructive completeness and the order of choices}

Start with a finite $\Ccal$ refutation. The choices must be made in the
following order; this prevents a circular size or congruence argument.
\begin{enumerate}[leftmargin=*,itemsep=0.35em]
\item Choose a cofinal horizon as in Theorem~\ref{g210:thm:cofinal}, exceeding
the refutation length and all finite local parameters; pad by truth steps.
Choose positive $F$ and its standalone Pell witnesses with the odd binomial
selector property. This fixes $B=\binom{F+t}{t}$.
\item Form the finite weighted pending-state path. Choose the odd multiplier
of $F$ so $y$ bounds this path and every prime at most $t$ divides $y$.
Finite CRT chooses an arbitrarily large positive beta code $x$ satisfying the
state residues and endpoints. Put $N=B^x$.
\item The generalized binomial CRT supplies the nine representative classes
modulo $N^2$. Choose the eight ordinary representatives positive first.
Because $N$ is odd, the ninth class has cofinally many even representatives.
Choose $C$ even and larger than every other nominal index and representative.
The prescribed index class gives $B\mid C+F+1$.
\item Set $D=2(C+F+1)$ and $q=D^x$ only now. Then $N\mid q$ and $D>2C$.
Evaluate the true requested power outputs, then choose the main square frame
and the positive canonical main Pell pairs after all requested atoms have
been fixed. The inherited cofinal anchor
completion supplies the private comparison witnesses, orientation quotient
and cap slack with one fewer principal pair.
\item Take the normalized low remainders. The estimates of
Section~\ref{g210:sec:radix} satisfy the retained common norm budget. Every
nonnegative slack has the inherited representation
$r^2+s^2+w(w+1)$, supplying the remaining integer witnesses.
\end{enumerate}
This constructs a zero of the changed arithmetic source from every finite
refutation. No guessed universal-machine parameter or unknown constant is
required.

\subsection{The precise logical dependency}
Combining the two directions with the square-shell numerical-calculus
calibration of Section~\ref{sec:foundation} proves the same syntactic
ZFC-refutation equivalence. The model-theoretic bridge is internal to the
background theory; it does not assume an externally well-founded transitive
model. Selected logical sources are retained, but no complete transitive
Lean or Rocq build is claimed.
